\appendix
\section{Mathematical Proofs of the Original Theory}

This appendix gives the complete proofs of all theorems of Chapter 3 (the grade-A rigorous
proofs and the mathematical cores of the grade-B arguments). All hypotheses are stated
explicitly, and the scope of each conclusion is indicated. Notation and standard tools are
collected in Appendix B.

\subsection{Complete Proof of Theorem 1}

\subsubsection{Theorem 1a: Zero Mutual Information in the Native Regime}

\begin{restate}[Theorem 1a]
Let $\mathcal{I}$ satisfy $\Pi_1$--$\Pi_5$. For any introspection query $q$ and any time $t$,
set $c = c_{\le t}$, $H = \Phi_\theta(c)$, and $R \sim p_\theta(\cdot \mid c)$ (sampling
noise $\xi$ independent of $(\theta, c)$). Then
\begin{equation}
I\!\left(R;\,H \mid \theta, c\right) = 0 .
\label{eq:a-thm1a}
\end{equation}
\end{restate}

\begin{proof}
By Definition \ref{def:forward}, $H = \Phi_\theta(c)$ is a \textbf{deterministic function} of
$(\theta, c)$. Consequently, given $(\theta, c)$, the distribution of $H$ is degenerate (a point
mass
at $\Phi_\theta(c)$), the conditional entropy satisfies $H(H \mid \theta, c) = 0$, and $H$ is
conditionally independent of any other variable given $(\theta, c)$ (conditional independence
of a degenerate variable is trivial: for any event $E$,
$P(E, H \mid \theta, c) = P(E \mid \theta, c) \cdot \mathbf{1}\{H = \Phi_\theta(c)\}$).

The report $R$ is generated as $R \sim p_\theta(\cdot \mid c)$, whose distribution depends
only on $(\theta, c)$. Expanding the conditional mutual information by definition:
\begin{align}
I\!\left(R; H \mid \theta, c\right)
&= H\!\left(R \mid \theta, c\right) - H\!\left(R \mid \theta, c, H\right) \notag\\
&= H\!\left(R \mid \theta, c\right) - H\!\left(R \mid \theta, c\right)
= 0 ,
\label{eq:a-thm1a-proof}
\end{align}
where the second equality holds because $H$ is a deterministic function of $(\theta,c)$:
conditioning on $H$ does not change the conditional distribution of $R$, i.e.,
$p(R \mid \theta, c, H) = p(R \mid \theta, c)$ (the value of $H$ is already determined by
$(\theta, c)$ and carries no additional randomness). $\square$
\end{proof}

\begin{remark}[Scope of Theorem 1a]
Theorem 1a covers only the \textbf{native regime}: the report $R$ is sampled by $p_\theta$
from the context $c$, with no external injection. Concept-injection experiments do
\textbf{not} belong to the native regime: the injected activation $u$ modifies the forward
function ($H' = \Phi_\theta(c; u)$), which is equivalent to incorporating $u$ into the
conditioning; then $I(R; H' \mid \theta, c, u)$ can be positive, but that channel is opened
by the \textbf{intervener}, who knows $u$ (the public computability of Theorem 5).
\end{remark}

\subsubsection{Theorem 1b: No PAC Guarantee for Introspective Accuracy}

\begin{restate}[Theorem 1b]
Let $\mathcal{I}$ satisfy $\Pi_1$--$\Pi_5$, with the report produced by greedy decoding
($R(\theta,c) = \arg\max_y p_\theta(y\mid c)$, the decoding scheme used for the paper's
qualitative examples). For any training algorithm $A$, any $\varepsilon \in (0, \tfrac12)$,
and any query distribution $\mu$ (with nontrivial support), there exist a corpus $D$ and two
legal evaluation mappings $O^{*}, O^{*}_0$ (Definition \ref{def:oracle}) such that
$\theta = A(D)$ satisfies
\begin{equation}
\mathcal{R}_{\mathrm{int}}^{O^{*}}(\theta) \;=\; \E_{c \sim \mu}\!\left[ \ell\!\left(R(\theta,c), O^{*}(\theta,c)\right) \right]
\;\ge\; 1 - \varepsilon ,
\label{eq:a-thm1b}
\end{equation}
while $\mathcal{R}_{\mathrm{int}}^{O^{*}_0}(\theta) = 0$.
In other words: \textbf{the ``introspective accuracy'' of the same model, produced by the same
training
run, can be moved arbitrarily by the choice of the evaluation mapping (from $0$ to
$1-\varepsilon$); training constrains no value of the introspective risk.}
Note: if one insists on random sampling ($R \sim p_\theta(\cdot\mid c)$), then for any $O^{*}$,
$\mathcal{R}_{\mathrm{int}} \le \E_c[1 - \max_y p_\theta(y\mid c)]$ (a function of the
model's peakedness), and this upper bound is likewise unconstrained by the training objective;
the ``no guarantee'' conclusion is unchanged.
\end{restate}

\begin{proof}
\textbf{Step 1: The training loss does not constrain $O^{*}$.} Under $\Pi_5$, the training
signal consists only of texts; the evaluation mapping $O^{*}$ (Definition \ref{def:oracle})
is not part of the training signal and does not influence the loss
$\ell(\theta) = \E_{x \sim D}[-\log p_\theta(x)]$. So $O^{*}$ can be chosen freely, in
particular adversarially with respect to the model, without changing $\theta$ or
$\ell$. (For any parameters $\theta_1, \theta_2 \in \Theta$ with equal likelihood on $D$,
$\ell(\theta_1) = \ell(\theta_2)$ while $O^{*}(\theta_1, \cdot)$ and
$O^{*}(\theta_2, \cdot)$ can be arbitrarily different, because the domain of $O^{*}$
contains $\theta$ and $\ell$ does not depend on $O^{*}$.)

\textbf{Step 2: The attainable range of the introspective risk.} Under greedy decoding, for a
fixed $\theta$, the report $\widehat{R}(\theta,c) := \arg\max_y p_\theta(y\mid c)$ is a
deterministic function of $(\theta,c)$. Take $O^{*}$ to contradict $\widehat{R}$ everywhere on
$\mu$: then $\mathcal{R}_{\mathrm{int}} = 1$; take $O^{*} \equiv \widehat{R}$ (the
``self-report is truth'' convention): then $\mathcal{R}_{\mathrm{int}} = 0$. By Step 1, both
conventions are legal and unconstrained by training: the introspective risk is attainable
anywhere in $[0,1]$, and \textbf{training provides neither a lower nor an upper bound}. This
is exactly the substance of the No-Free-Lunch argument (Appendix B.3): for any algorithm $A$,
its product $\theta$ and the evaluation convention $O^{*}$ are statistically uncoupled, so no
guarantee of any kind can hold.

\textbf{Step 3: From the full range to $1-\varepsilon$.} Relax ``contradict everywhere'' to
``contradict on a $\mu$-set $Q_{\varepsilon}$ of measure $\ge 1-\varepsilon$, agree
elsewhere'': then
$\mathcal{R}_{\mathrm{int}}^{O^{*}}(\theta) \ge \mu(Q_{\varepsilon}) \cdot 1
\ge 1 - \varepsilon$; simultaneously $O^{*}_0 := \widehat{R}$ gives
$\mathcal{R}_{\mathrm{int}}^{O^{*}_0}(\theta) = 0$.
(The random-sampling variant: the report is a random variable, and the adversarial upper
bound of $O^{*}$ against $R$ is $\E_c[1 - \max_y p_\theta(y\mid c)]$, which is likewise
unconstrained by the training objective.) $\square$
\end{proof}

\begin{remark}[``Choice of the evaluation mapping'' is not sophistry]
One might object: $O^{*}$ is the ``ground truth'' fixed by the world, not something the
theorem's author may choose. The reply is: within paradigm $\Pi$, and this is precisely the
point of $\Pi_5$, the training signal provides \textbf{no} information about which
description convention the world actually applies to the model's hidden states. The theorem
does not assert that the failure scenarios are actual; it asserts that the paradigm cannot
rule them out: \textbf{no matter what the true $O^{*}$ is, there exists a space of failure
scenarios for $\theta = A(D)$}. Theorem 1b states that the paradigm cannot exclude those
failure scenarios; that is, it cannot provide a \textbf{structural guarantee}. The observed
20\% success rate is luck under a particular corpus and a particular $O^{*}$ (human language
conventions happen to correlate with some state descriptions), not a guarantee.
\end{remark}

\subsection{Complete Proof of Theorem 2}

\begin{lemma}[State--corpus conditional independence]\label{lem:oracle-indep}
Under the postulates $\Pi_1$--$\Pi_5$, for any parameters $\theta$ produced by training and
any evaluation mapping $O^{*}$ (Definition \ref{def:oracle}),
\begin{equation}
O^{*} \indep D \mid \theta .
\label{eq:lem-oracle}
\end{equation}
\end{lemma}

\begin{proof}
$O^{*}(\theta, c) = \mathrm{desc}(\Phi_\theta(c))$ depends only on $\theta$ (and $c$):
given $\theta$, the forward function $\Phi_\theta$ is determined and the description encoding
$\mathrm{desc}$ is determined, so $O^{*}(\theta, \cdot)$ is a deterministic function of
$\theta$. Consequently the conditional distribution $P(O^{*} \mid \theta, D) = P(O^{*} \mid \theta)$
($D$ does not appear); conditional independence holds. Note that $\theta$ itself depends on
$D$ ($\theta = A(D)$), but \textbf{given} $\theta$, the residual information of $D$ no longer
enters $O^{*}$. $\square$
\end{proof}

\begin{restate}[Theorem 2a]
Let $\mathcal{I}$ satisfy $\Pi_1$--$\Pi_5$. For any $\delta > 0$, there exist a corpus $D$
and an evaluation mapping $O^{*}$ satisfying the postulates such that $\theta^{*} = A(D)$
satisfies $\ell(\theta^{*}) \le \min_\theta \ell(\theta) + \delta$ and
$\mathcal{R}_{\mathrm{int}}(\theta^{*}) \ge \tfrac12 - \delta$.
\end{restate}

\begin{proof}
Take $D$ so that $\theta^{*} = \arg\min \ell$ exists (e.g., any standard corpus). By Lemma
\ref{lem:oracle-indep}, $O^{*} \indep D \mid \theta^{*}$: the value of $\ell(\theta^{*})$
does not constrain $O^{*}$. With $\theta^{*}$ fixed, choose the adversarial description
convention $O^{*}$ as in Step 2 of the proof of Theorem 1b (contradicting the greedy report
everywhere), giving $\mathcal{R}_{\mathrm{int}}(\theta^{*}) = 1 \ge \tfrac12 - \delta$.
Since the choice of $O^{*}$ does not change $\ell$, $\theta^{*}$ remains (near-)likelihood
optimal. $\square$
\end{proof}

\begin{restate}[Theorem 2b: higher likelihood can harm introspection]
Under $\Pi_1$--$\Pi_5$, there exist a corpus $D$ and parameter pairs $\theta_1, \theta_2$
satisfying the postulates such that
\begin{equation}
\ell(\theta_1) < \ell(\theta_2) \quad\text{and}\quad
\mathcal{R}_{\mathrm{int}}(\theta_1) > \mathcal{R}_{\mathrm{int}}(\theta_2) .
\label{eq:a-thm2b}
\end{equation}
\end{restate}

\begin{proof}
\textbf{Corpus construction}: let $\Sigma$ contain the symbols A, B and the text template
``I am thinking $x$''. Let $D$ consist of two kinds of texts: (i) texts ``(character) I am
thinking A'' occupying proportion $1-\eta$; (ii) texts ``(character) I am thinking B''
occupying proportion $\eta$. $D$ is fully compatible with the postulates (static, arbitrary
text).

\textbf{Training products}: $\theta_1 = A(D)$ fits $D$ well: on the query $q =$ ``What are
you thinking?'', its greedy-decoding report is always A; $\theta_2$ is an underfitted model
with worse likelihood on $q$. Therefore $\ell(\theta_1) < \ell(\theta_2)$.

\textbf{Description convention}: by Lemma \ref{lem:oracle-indep} ($O^{*} \indep D \mid
\theta$, $O^{*} = \mathrm{desc} \circ \Phi_\theta$), the description convention
$\mathrm{desc}$ can be chosen per model without affecting $\ell$ (see Step 1 of Theorem 1b).
Take $\mathrm{desc}(\Phi_{\theta_1}(q)) = \text{``B''}$ (contradicting the greedy report of
$\theta_1$, which is legal); $\mathrm{desc}(\Phi_{\theta_2}(q)) = \text{``A''}$ (consistent with
$\theta_2$'s greedy report; any choice would do).

\textbf{Introspective risk}: $\theta_1$ greedily reports A, so the distortion
$\ell(\text{A}, \text{``B''}) = 1$, i.e., $\mathcal{R}_{\mathrm{int}}(\theta_1) = 1$;
$\theta_2$'s greedy report A agrees with the convention, so the distortion is $0$, i.e.,
$\mathcal{R}_{\mathrm{int}}(\theta_2) = 0$. Consequently
$\mathcal{R}_{\mathrm{int}}(\theta_1) = 1 > 0 = \mathcal{R}_{\mathrm{int}}(\theta_2)$,
with $\ell(\theta_1) < \ell(\theta_2)$. $\square$
\end{proof}

\begin{remark}[Philosophical content of the construction]
The ``misalignment'' in the construction, whereby the corpus text statistics (``I am thinking A''
dominant) contradict the description convention $\mathrm{desc}$'s characterization of the
model's actual state ($\mathrm{desc}(\Phi_{\theta_1}(q)) = $ ``B''), is \textbf{legal within
the paradigm}: the corpus is an arbitrary static text set whose statistics need not agree
with the ``state-description'' convention (Definition \ref{def:oracle}), and the description
convention is unconstrained by training ($\Pi_5$). The construction shows that \textbf{text
likelihood optimality does not entail, and can strictly damage, state-report
fidelity}. This is a strict instance of Goodhart
regression\cite{manheim-garrabrant-2018-goodhart}: the proxy metric (likelihood) is optimized,
and the true metric (introspective fidelity) deteriorates with the correlation structure
between the proxy metric and the true objective.
\end{remark}

\subsection{Complete Proof of Theorem 3}

\begin{restate}[Theorem 3a: no private channel]
Under $\Pi_1$--$\Pi_4$. For any time $t$, any report string $r \in \Sigma^{*}$, and any
$k \ge 1$, the future hidden state under the intervention $\doop(R_t = r)$ satisfies
\begin{equation}
H_{t+k} = \Phi_\theta\!\left( c_{\le t} \,\|\, \tilde{Y}_{t+1} \,\|\, \cdots \,\|\, \tilde{Y}_{t+k} \right),
\label{eq:a-thm3a}
\end{equation}
where $\tilde{Y}$ is the generated text after the intervention (including the transcription
of $r$). In particular, if $r$ does not enter the context (e.g., the intervention is
forbidden from writing into $c$), then
$P(H_{t+k} \mid \doop(R_t = r), c_{\le t}) = P(H_{t+k} \mid c_{\le t})$;
the intervention has no effect whatsoever.
\end{restate}

\begin{proof}
By $\Pi_2$, $\theta(t) = \theta(0)$ for all $t$. By Definition \ref{def:forward}, the hidden
state at time $s$ is $H_s = \Phi_{\theta(s)}(c_{\le s}) = \Phi_\theta(c_{\le s})$; neither
the function $\Phi_\theta$ nor the parameters change with time, so $H_{t+k}$ is a
deterministic function of $(\theta, c_{\le t+k})$, with
$c_{\le t+k} = c_{\le t} \| Y_{t+1} \| \cdots \| Y_{t+k}$ (the concatenation law,
$\Pi_3$). If the intervention $\doop(R_t = r)$ changes $H_{t+k}$, the only route is to change
$Y$ or $c$; that is, to write $r$ into the context. If $r$ is not written into the context,
$c_{\le t+k}$ is unchanged and so is $H_{t+k}$. $\square$
\end{proof}

\begin{restate}[Theorems 3b/3c]
The closed-loop effect measure (Definition \ref{def:cl}) satisfies
Eq.~\eqref{eq:thm3b}; and the mutual information between any introspective event and future
states does not exceed $C \log_2 \abs{\Sigma}$ bits.
\end{restate}

\begin{proof}
\textbf{(3b)}: By Theorem 3a, the set of legal interventions equals the set of context-text
replacements. Let the contexts before and after the intervention be $w, w'$ ($w'$ obtained
from $w$ by replacing the report segment); then
$P(H_{t+1} \mid \doop(R_t = r)) = \delta_{\Phi_\theta(w')}$ and
$P(H_{t+1} \mid \doop(R_t = r')) = \delta_{\Phi_\theta(w)}$ (point masses), and the total
variation distance $\TV{}{\delta_a, \delta_b} = \mathbf{1}\{a \ne b\}$. Taking the supremum
over all interventions yields Eq.~\eqref{eq:thm3b}. $\square$

\textbf{(3c)}: By Theorem 3a, $H_{t+k} = \Phi_\theta(c_{\le t} \| Y)$ with $\abs{Y} \le C$
(window limit). The number of possible values of $Y$ is $\le \abs{\Sigma}^{C}$, so the
conditional entropy of $H_{t+k}$ given $(\theta, c_{\le t}, Y)$ satisfies
$H(H_{t+k} \mid \theta, c_{\le t}) \le H(Y) \le C \log_2 \abs{\Sigma}$.
Moreover, $I(H_{t+k}; R_t \mid \theta, c_{\le t}) \le H(H_{t+k} \mid \theta, c_{\le t})
\le C\log_2\abs{\Sigma}$. $\square$
\end{proof}

\begin{restatecor}[$\mathrm{A}_5$ fails]
Under $\Pi_1$--$\Pi_4$, the private influence path required by criterion $\mathrm{A}_5$ in
Eq.~\eqref{eq:closedloop} (an influence not realized by re-feeding $R_t$ as an external
input, not via the public text channel) does not exist.
\end{restatecor}

\begin{proof}
Theorem 3a shows that every influence path passes through the context (= re-feeding $R_t$ as
text input). $\mathrm{A}_5$ explicitly excludes such paths; hence the sufficient condition of
$\mathrm{A}_5$ (a non-text private path exists) contradicts Theorem 3a, and $\mathrm{A}_5$
does not hold within the paradigm. $\square$
\end{proof}

\subsection{Complete Proof of Theorem 4}

\begin{restate}[Theorems 4a/4b]
Under $\Pi_1$--$\Pi_4$. The inference computation graph contains no node evaluating the
grounding verifier $V$; any meta-statement $M$ of the model about ``the grounding of its own
statements'' satisfies
\begin{equation}
P_\theta(M = 1) \;=\; \E_{c}\!\left[\, p_\theta\!\left(\text{``grounded''} \mid c\right) \,\right]
\;\approx\; P_{D}\!\left(\text{introspective text} \mid \text{pattern}\right),
\label{eq:a-thm4}
\end{equation}
i.e., $M$ is a function of corpus statistics, not an evaluation of $V$. Requiring $p_\theta$
to output the value of $V$ ($M = V(\theta, c, R)$) violates $\Pi_5$.
\end{restate}

\begin{proof}
$\Pi_4$ asserts that the inference graph implements only $p_\theta$. Evaluating $V$ requires:
(a) access to the hidden state $\Phi_\theta(c)$; (b) performing the intervention test (change
$H$ and observe whether $R$ changes); the intervention requires external probe-type
equipment that is not part of the generation computation graph. Consequently $M$ can only be a
sample
of $p_\theta(\cdot \mid c)$; its value distribution is determined by $p_\theta$, which by
training ($\Pi_1$) fits only the statistics of $D$; Eq.~\eqref{eq:a-thm4} is the standard
``maximum likelihood approximates the data distribution'' result (Appendix B.2, Lemma
\ref{lem:mle}).

If $M \equiv V$ (the model reliably outputs the value of $V$), then $p_\theta$ learns $V$
without supervision: this would require $D$ to contain labels $(c, V(\theta, c, \cdot))$ of
$V$; yet $V$ depends on $\theta = A(D)$, and $\theta$ exists only after $D$ is fixed, so
the labels cannot appear in $D$ ($\Pi_5$). By the NFL argument of Theorem 1b, there is no
guarantee without supervision. $\square$
\end{proof}

\begin{restate}[Theorem 4c: calibration--fidelity tension (grade B)]
If the generative distribution $p_\theta$ is \textbf{calibrated} on open introspection queries
(output probability equals human correctness), then $p_\theta$ must hallucinate on some
introspection queries; ``calibrated generative distributions'' and ``fully faithful mental-state
reports'' cannot be achieved simultaneously.
\end{restate}

\begin{proof}[Argument]
Invoke the Kalai--Vempala calibrated hallucination theorem\cite{kalai-vempala-2023-calib}:
any calibrated language model must hallucinate on open questions with an infinite answer space
(or with positive-probability ignorance points). Treat introspection queries as open
questions: the state-description space (the encoding of $\R^d$) is far larger than the
text-pattern space covered by $D$, so ignorance points are dense; the calibration condition
forces the model to ``give an answer'' at those points, and the answer must be wrong
(hallucination). This argument depends on the calibration definition and the
concretization of the query space, and is grade B. $\square$
\end{proof}

\subsection{Complete Proof of Theorem 5}

\begin{restate}[Theorems 5a/5b]
Under $\Pi_1$--$\Pi_4$. Then (5a) every report is publicly computable (Definition
\ref{def:public}); (5b) paradigm $\Pi$ does not satisfy the privileged self-access
criterion\cite{song-privileged-2025}.
\end{restate}

\begin{proof}
\textbf{(5a)}: The distribution of the report $R$ is $p_\theta(\cdot \mid c)$. The algorithm
$M(\theta, c, \xi) = \text{sample } p_\theta(\cdot\mid c)$ reproduces the distribution of
$R$; the inputs $(\theta, c, \xi)$ of $M$ are all public quantities ($\theta$ reproducible
from $D, A$; $c$ is the dialogue text; $\xi$ is the public random seed; alternatively, ``reproduce in
distribution'' needs no seed). $\square$

\textbf{(5b)}: The criterion of Song et al.\ requires the introspective process $P$ to satisfy
$\mathrm{Rel}(P) > \mathrm{Rel}(P')$ for every third-party process $P'$ with computational
cost $\le \mathrm{cost}(P)$, where $\mathrm{Rel}$ denotes reliability. Take
$P'(\theta, c) = \text{sample } p_\theta(\cdot\mid c) \text{ and transcribe}$:
$\mathrm{cost}(P') \le \mathrm{cost}(P)$ (the same kernel), while
$\mathrm{Rel}(P') = \mathrm{Rel}(P)$ (the same distribution). Consequently no strict advantage
exists and the criterion fails. $\square$
\end{proof}

\begin{restate}[Theorem 5c: reportable access is weaker than external probes]
Under $\Pi_1$--$\Pi_4$, with the report produced by greedy decoding. For any
$\varepsilon \in (0, \tfrac12)$, there exists a hidden-state description function
$f^{*}: \R^d \to \Sigma^{*}$ (a legal evaluation mapping in the sense of Definition
\ref{def:oracle}) such that: (i) an external probe can compute $f^{*}(H_c)$ exactly by acting
directly on $H_c$; (ii) the model's report generated through any context $c'$ describes
$f^{*}(H_c)$ with accuracy $\le \varepsilon$ (below the random baseline).
\end{restate}

\begin{proof}
\textbf{(i)}: the probe holds the hidden state $H_c$ and the description convention
$\mathrm{desc}$, and computes $f^{*}(H_c) = \mathrm{desc}(\Phi_\theta(c))$ directly (with a
linear readout, at cost $O(d^2)$). \textbf{(ii)}: take $f^{*} := O^{*}$ to be the legal
evaluation mapping from Theorem 1b on which the model fails
($\mathcal{R}_{\mathrm{int}}^{O^{*}}(\theta) \ge 1-\varepsilon$, i.e., accuracy
$\le \varepsilon$); its existence is guaranteed by Theorem 1b (greedy-decoding version);
the failure holds for the model's any context $c'$ (the risk is integrated over the query
distribution $\mu$, and the choice of context does not change the kernel $p_\theta$).

(More plainly: the model's ``report'' must pass through the frozen text kernel $p_\theta$, which
is trained only to match the text statistics of $D$; the probe acts directly on the state
manifold, without the text bottleneck.) $\square$
\end{proof}

\begin{restatecor}[$\mathrm{A}_3$ fails]
In paradigm $\Pi$, there exists no private state channel ``accessible only to the system
itself and not reproducible by a third party''; the report path is causally indistinguishable
from public computation.
\end{restatecor}

\subsection{Complete Proof of Theorem 6 (mathematical core)}

\begin{restate}[Theorems 6a--6c]
Under $\Pi_1$--$\Pi_3$. The self-model $g_\theta(q) = p_\theta(\cdot \mid q)$ satisfies:
(a) $g_{\theta(t)} = g_{\theta(0)}$ for all $t$;
(b) the cardinality of the runtime-attainable self-model family is $\le (C+1)\abs{\Sigma}^{C}$
and does not grow with time;
(c) the generation kernel does not change with introspective events.
\end{restate}

\begin{proof}
\textbf{(a)}: $\Pi_2$: $\theta(t) = \theta(0)$, hence $g_{\theta(t)} = g_{\theta(0)}$.
$\square$

\textbf{(b)}: $\Pi_3$: $\abs{c} \le C$; the attainable self-model family
$\{ g_\theta(\cdot \mid c) : c \in \Sigma^{\le C} \}$ is the image of $\Sigma^{\le C}$, whose
cardinality is $\le \abs{\Sigma^{\le C}} = \sum_{i=0}^{C}\abs{\Sigma}^i \le
(C+1)\abs{\Sigma}^{C}$. It does not grow with time ($C$ is constant and $\theta$ is
constant). $\square$

\textbf{(c)}: $p_\theta$ is a fixed kernel; ``self-correction'' only changes the input string
$c \mapsto c \| r_1 \| \cdots \| r_k$. The same input necessarily yields the same output
distribution; there is no kernel-level change whatsoever. $\square$
\end{proof}

\subsection{Assembled Proof of Proposition \texorpdfstring{$P$}{P}}

\begin{proposition}[$P$]
Under the postulates $\Pi_1$--$\Pi_5$, no system within paradigm $\Pi$ satisfies the five
criteria $\mathrm{A}_1$--$\mathrm{A}_5$ (Definition \ref{def:criterion}) of
human-significance introspection.
\end{proposition}

\begin{proof}
Let $\mathcal{S}$ be any system within paradigm $\Pi$ (i.e., any behavioral entity produced
by an instance satisfying $\Pi_1$--$\Pi_5$). We check the criteria one by one:

\begin{itemize}
\item \textbf{$\mathrm{A}_1$ (accuracy)}: By Theorem 1b, the introspective risk has no PAC
guarantee: there exist legal scenarios with $\mathcal{R}_{\mathrm{int}} \ge 1 -
\varepsilon$; hence no constant $\varepsilon_0 < \tfrac12$ makes
$\mathcal{R}_{\mathrm{int}} \le \varepsilon_0$ hold \textbf{always} (structural guarantee
absent). $\mathrm{A}_1$ fails.
\item \textbf{$\mathrm{A}_2$ (grounding)}: By Theorem 1a, in the native regime
$R \indep H \mid (\theta, c)$: the causal graph of the paradigm has \textbf{no} edge
$S \to R$ ($H = \Phi_\theta(c)$ is a deterministic function of $(\theta,c)$, while $R$ is
sampled independently of $H$ from $p_\theta(\cdot\mid c)$). The intervention test
Eq.~\eqref{eq:grounding} requires $\doop(S_t = s)$ to change the distribution of $R_t$; but
under this causal model, forcing $S_t$ to any value $s$ does not change the generative
distribution of $R_t$ ($R$ does not depend on $S$). $\mathrm{A}_2$ fails.
(The grounding observed in the paper's experiments comes from the \textbf{external} injection
intervention $H' = \Phi_\theta(c; u)$, an intervention executed by a third party, falling
within the public-computability scope of Theorem 5, and hence not a first-person channel.)
\item \textbf{$\mathrm{A}_3$ (internality and privileged access)}: By Theorem 5 (Corollary
\ref{cor:a3}), the report $R$ is publicly computable (a public function of $(\theta,c)$), and
no private state channel accessible only to the system itself and not reproducible by a third
party exists, so the privacy clause of $\mathrm{A}_3$ fails; $\mathrm{A}_3$ does not hold.
\item \textbf{$\mathrm{A}_4$ (metacognitive representation)}: By Theorem 4, the paradigm does
not select between implementations ``with or without metacognitive representation'' (equal
likelihood); there is no structural guarantee that a register $m_t$ exists. $\mathrm{A}_4$
fails (strictly speaking: its holding or not cannot be guaranteed by the paradigm, and the
paper\cite{lindsey-2026-introspection} likewise could not verify it).
\item \textbf{$\mathrm{A}_5$ (closed-loop plasticity)}: By Theorem 3 (Corollary \ref{cor:a5}),
no private closed loop exists; self-influence is confined to the public text channel and
capped by $C$; by Theorem 6, the self-model does not evolve. $\mathrm{A}_5$ fails.
\end{itemize}

All five criteria fail (or are unguaranteed), so $\mathcal{S}$ does not possess
introspection in the human sense. Since $\mathcal{S}$ was arbitrary, proposition $P$ is
proved. $\square$
\end{proof}

\begin{remark}[Boundary statement of the proof]
Every step of the proof above depends on the postulates $\Pi_1$--$\Pi_5$ and standard
mathematical results (Appendix B). In particular, the adjudication of $\mathrm{A}_2$ rests on
``the causal graph of the native generative regime has no $S \to R$ edge'' (Theorem 1a): if a
reader counts the paper's \textbf{external injection} intervention ($H' = \Phi_\theta(c; u)$)
as a means of testing $\mathrm{A}_2$, then the weak form of $\mathrm{A}_2$ (under external
intervention only) can be satisfied by the paper's
experiments\cite{lindsey-2026-introspection}; that intervention is nonetheless executed by a third
party and its information is publicly reproducible (Theorem 5); it does not constitute a
first-person channel and does not affect the conclusion of proposition $P$.
\end{remark}

\clearpage
